\section{Discussion}
\label{sec:discussion}

\textbf{Existence and authority are different checks.} A verifier over an index of names and citations catches fabricated cases, and feedback from it or search over it removed almost all of them. It did not help the models find the judgment that decides a question, and name search made that rarer, because the models searched with the question's words and listed any judgment whose parties happened to share them. A real but irrelevant judgment passes an existence check. A system that shows only ``verified'' citations can therefore look more reliable while citing worse authorities. Finding the authority needs what the metadata lacks: a description of what each judgment decided, which the reporter's headnote supplies, and a record of which judgments the Court itself relies on, which the citation graph supplies. With both, the same small models named a deciding judgment in \RefAllTopicPrior{} of \RefQuestionsAll{} cases and almost never cited a case that does not exist.

\textbf{The retriever does the work.} Listing the retriever's top three results answers about as many questions as the models that use it (\QHybridPriorHitThree{} against \RefTopicPriorGemma--\RefTopicPriorLlama), and the models listed almost only what the search had shown them. For a system that must not invent authorities this is the desirable division of labour: the index guarantees existence, the retriever and the graph supply relevance, and the model at most phrases and orders the result. The citation prior helps because it encodes the Court's own judgment of importance, as network measures of precedent do for the United States Supreme Court \cite{fowler2007}. It favours old and frequently cited judgments, including judgments later overruled, such as \emph{A.D.M.\ Jabalpur}, which a lawyer must cite as overruled. Popularity alone is useless (recall@10 \RecPopRTen), so the prior can only reorder judgments that already match the subject.

\textbf{Declining one question is not reliability.} The two Qwen models declined every request to identify a judgment from its citation, which looks like calibrated caution. Yet asked, with the same system prompt, for the citation of a judgment they had just been told exists, they produced one for most judgments, and nearly all were wrong. Whether a small model admits ignorance therefore depends on how the question is framed, and a single abstention test says little about how it will behave in drafting, where it is asked to supply citations rather than check them.

\textbf{Neutral citations invite invention.} The Supreme Court began issuing neutral citations in 2023 and assigned them retrospectively to its reported judgments from 1950 \cite{scneutral,scneutral2}, so for most judgments they are recent additions that a model is unlikely to have seen in use. The models nevertheless produced them readily, many with numbers that were never issued. The same property that makes them easy to fabricate makes them easy to check: for the \INSCCompleteCount{} years whose numbering the index holds completely, a single lookup settles whether a neutral citation exists.

\textbf{What remains unchecked.} Neither the verifier nor the retriever checks that a judgment supports the proposition it is cited for, which needs the judgment's reasoning and is hard even for strong models \cite{verma2026}. The graph records that one judgment cites another, not whether it follows, distinguishes or overrules it. Citations to the SCC and AIR reports are resolved through parallel S.C.R. citations or by party names, never by their own numbers, and judgments reported only in those reports are not in the index: \UnresNotIndex{} of the \UnresN{} unresolved citations we inspected fall in that gap. Recent neutral citations can only be confirmed, not refuted, because the index holds only reported judgments after 2013.

\textbf{Limitations.} We evaluated four open-weight models of roughly 4 to 8 billion parameters (Table~\ref{tab:models}) in 4-bit quantisation, with one prompt and greedy decoding. Larger and commercial systems may know many more landmark citations, and different prompts may change the balance between answering and declining. The benchmark is small (\NumJudgments{} judgments, \NumTopics{} questions) and the landmark set is our own selection. The audit labels, the reference answers and the manual check of the graph come from a single annotator, working from the index, the judgment text and published sources, so we cannot report inter-annotator agreement. We compiled the reference answers after the closed-book runs. They list only judgments that decided each question, so they measure relevance only from below, and because they are landmarks they favour the citation prior. The recommendation experiment has neither problem, but it treats the judgments a decision cites as its relevant authorities, which is incomplete: a court cites some of the relevant judgments, not all of them. Each search condition used one tool description and one query method. The frozen name parser misreads a name that ends in ``v.''\ with no respondent (two answers affected) and a party name that begins with the initial V. VERIFIED means that a judgment with the named parties exists near the stated year, not that it is the authority the model meant. The corpus is a snapshot of an archive that is updated periodically, and it has its own errors: \SCDupGroups{} neutral citations are attached to more than one record, titles contain OCR errors, some judgments are filed under the wrong year, and \GraphTextFail{} PDFs are missing.

\textbf{Serving cost.} Everything runs on one laptop with default PostgreSQL settings. We timed \LatQueries{} queries to the index over HTTP, each run twice, directly after a database restart: exact lookups by CNR, neutral or S.C.R. citation took a median of \LatExactMedian{}\,ms (at most \LatExactMax{}\,ms), and keyword queries for common phrases a median of \LatCommonFirstMedian{}\,ms on first touch and \LatCommonWarmMedian{}\,ms when repeated. The verifier checked a cited case in a median of \VerifyMedianMs{}\,ms (95th percentile \VerifyPNinetyFiveMs{}\,ms). Embedding the headnotes ran at about \EmbedRate{} judgments per second on the laptop GPU, and a retrieval query, including embedding the query, took about \QueryMs{}\,ms on average.
